\documentclass{solve}

\usepackage{tikz}

\begin{document}
  \pageHeader{Odwrotność}{ODW0}{64}{1}

  Jasio często lubi robić rzeczy na~opak. W~dodatku jest dość leniwy. Tym razem
  na opak oraz leniwie postanowił napisać napis. Za każdym razem kiedy Jasiowi
  przyszło do głowy coś odwrócić to umieszczał to w~nawiasie: na~przykład, jeśli
  Jasio napisał \texttt{inf(ormat)yka} to miał naprawdę na myśli napis
  \texttt{inftamroyka}. Problem w tym, że Jasio jest bardzo sprytny i czasami
  nawiasy umieszczał wewnątrz innych nawiasów. Przy odczytywaniu prawdziwych
  intencji Jasia oczywiście należy (tak jak w~przypadku działań matematycznych)
  odwracać zawartość nawiasów od najbardziej do najmniej zagnieżdżonych:
  na~przykład: \texttt{nap(is(odw)ro(co)nypar)erazy} należy rozumieć jako
  równoważny z \texttt{nap(iswdoroocnypar)erazy}, czyli:
  \texttt{naprapyncoorodwsierazy}.

  Jasio oczywiście nie ma teraz chęci na ustalanie co~tak naprawdę miał
  na myśli. Ale Ty przecież potrafisz programować...

  Napisz program, który wczyta napis napisany przez Jasia, wyznaczy co tak
  naprawdę Jasio miał na myśli i~wypisze wynik na wyjście.

	\section{Wejście}
  W~pierwszym (jedynym) wierszu wejścia znajduje się ciąg małych liter alfabetu
  angielskiego oraz nawiasów \texttt{(} i \texttt{)} -- napis Jasia.

  Można założyć, że po zignorowaniu liter, wejście stanowi poprawne
  nawiasowanie.

	\section{Wyjście}
  W~pierwszym (jedynym) wierszu wyjścia należy wypisać napis Jasia po~odwróceniu
  fragmentów w~nawiasach, zgodnie z~regułami opisanymi powyżej.

  \section{Ograniczenia}
  Długość napisu Jasia na~wejściu nie~przekracza $1\,000\,000$ znaków.

	\section{Przykład}
    \testVerb{odw0.in}{odw0.out}{}
\end{document}
